\Askhsh{21197_SOLUTION}{1}
\begin{ylh}{1}
    \item [Ύλη από εικόνα]
\end{ylh}
\Askhseis{ΘΕΜΑ 4}
\begin{alist}
    \item[\textalpha)]
    \begin{rlist}
        \item Το ημικύκλιο με διάμετρο $\text{AB} = 2\alpha$ έχει ακτίνα $\alpha$ και εμβαδόν $\text{E}_{\text{AB}} = \frac{\pi\alpha^2}{2}$. Αφού το εμβαδόν του ημικυκλίου είναι 10 τότε:
        \[ \text{E}_{\text{AB}} = \frac{\pi\alpha^2}{2} \text{ ή } 10 = \frac{\pi\alpha^2}{2} \text{ ή } \pi\alpha^2 = 20 \text{ ή } \alpha^2 = \frac{20}{\pi} \]
        Το εμβαδόν του τετραγώνου $\text{ΑΒΓΔ}$ με πλευρά $2\alpha$ είναι:
        \[ (\text{ΑΒΓΔ}) = (2\alpha)^2 = 4\alpha^2 = 4 \cdot \frac{20}{\pi} = \frac{80}{\pi} \]

        \item Το σημείο $\Lambda$ είναι το μέσο της πλευράς $\text{ΓΔ}$ του τετραγώνου, επομένως $\Delta\Lambda = \frac{2\alpha}{2} = \alpha$.
        Στο ορθογώνιο τρίγωνο $\text{ΑΔΛ}$ εφαρμόζοντας το πυθαγόρειο θεώρημα έχουμε:
        \[ \text{A}\Lambda^2 = \text{A}\Delta^2 + \Delta\Lambda^2 = (2\alpha)^2 + \alpha^2 = 5\alpha^2 \]
        Από το α) i. ερώτημα είναι $\alpha^2 = \frac{20}{\pi}$, επομένως $\text{A}\Lambda^2 = 5\alpha^2 = 5 \cdot \frac{20}{\pi} = \frac{100}{\pi} \]
    \end{rlist}
    \item[\textbeta)]
    \begin{rlist}
        \item Το ζητούμενο εμβαδόν $\text{Ε}$ του σκιασμένου σχήματος, υπολογίζεται αν από το εμβαδόν του τεταρτοκυκλίου $\text{A}\widehat{\text{MN}}$ αφαιρέσουμε το εμβαδόν $\text{E}_{\text{AB}}$ του ημικυκλίου με διάμετρο την $\text{AB}$.
        Το εμβαδόν του τεταρτοκυκλίου $\text{A}\widehat{\text{MN}}$ είναι:
        \[ (\text{A}\widehat{\text{MN}}) = \frac{\pi \cdot \text{A}\Lambda^2}{4} = \frac{\pi \cdot \frac{100}{\pi}}{4} = \frac{100}{4} = 25 \]
        επομένως το ζητούμενο εμβαδόν $\text{Ε}$ του σκιασμένου σχήματος είναι:
        \[ \text{E} = (\text{A}\widehat{\text{MN}}) - \text{E}_{\text{AB}} = 25 - 10 = 15 \]

        \item Από το ερώτημα (β.i) το εμβαδόν του τεταρτοκυκλίου $\text{A}\widehat{\text{MN}}$ είναι $(\text{A}\widehat{\text{MN}}) = 25$ και από το α) i. ερώτημα το εμβαδόν του τετραγώνου $\text{ΑΒΓΔ}$ είναι $(\text{ΑΒΓΔ}) = \frac{80}{\pi}$, επομένως ο λόγος του εμβαδού του τεταρτοκυκλίου $(\text{A}\widehat{\text{MN}})$ προς το εμβαδόν του τετραγώνου $(\text{ΑΒΓΔ})$ θα είναι:
        \[ \frac{(\text{A}\widehat{\text{MN}})}{(\text{ΑΒΓΔ})} = \frac{25}{\frac{80}{\pi}} = \frac{25\pi}{80} = \frac{5\pi}{16} \]
    \end{rlist}
\end{alist}